% This is samplepaper.tex, a sample chapter demonstrating the
% LLNCS macro package for Springer Computer Science proceedings;
% Version 2.21 of 2022/01/12
%
\documentclass[runningheads]{llncs}
\usepackage{tikz}
\usepackage{subcaption}
\usepackage{url}
\usepackage{hyperref}


\usepackage{xcolor}
\usepackage{amsmath}
\usepackage{amsfonts}
\setcounter{secnumdepth}{3}
\DeclareMathOperator{\Truth}{truth}

%
\usepackage[T1]{fontenc}
\usepackage{graphicx}
\usepackage[labelfont=bf,labelsep=period]{caption}

\begin{document}
%

\title{A Cooperative and Context-Aware Approach for Employee Attrition Prevention}

\titlerunning{Cooperative, Context‑Aware Attrition Prevention}
% If the paper title is too long for the running head, you can set
% an abbreviated paper title here
%
\author{
Mariia Godz\inst{1}\orcidID{0009-0002-0183-7905} \and \\
Pavel Novoa-Hern\'andez\inst{2}\orcidID{0000-0003-3267-6753} \and \\
David A. Pelta\inst{1}\orcidID{0000-0002-7653-1452}
}

\authorrunning{M. Godz et al.}

\institute{
Dept. Computer Science and Artificial Intelligence,\\
University of Granada, Granada, Spain
\and
Dept. Computer Engineering and Systems,\\
University of La Laguna, La Laguna, Spain
}
%
\authorrunning{M. Godz et al.}
% First names are abbreviated in the running head.
% If there are more than two authors, 'et al.' is used.
%

%
\maketitle              % typeset the header of the contribution
%
\begin{abstract}
Employee attrition remains a persistent organizational challenge due to its economic cost, loss of accumulated knowledge, and impact on workforce stability.
While extensive research has focused on predicting attrition using machine learning, effective retention decisions require more than accurate risk estimation.
Organizations must prioritize limited intervention efforts by integrating heterogeneous assessments from different organizational units and accounting for contextual and policy-driven constraints.
This paper proposes a cooperative and context-aware decision approach for employee attrition intervention prioritization.
Multiple decision components, each reflecting a distinct organizational perspective and trained on different subsets of employee attributes, produce partial risk assessments.
These assessments are aggregated through confidence-weighted integration and modulated by factors that may influence context encoding organizational priorities.
Rather than competing with predictive performance benchmarks, this work demonstrates how heterogeneous automated assessments can be structured into a transparent and adaptable prioritization mechanism.
The approach is evaluated on a public employee attrition dataset as a methodological testbed for cooperative, context-aware decision-making under constrained resources.

\keywords{Employee attrition \and Cooperative decision-making \and Context-aware decision-making \and Automated decision systems \and Human resource analytics}
\end{abstract}



%
%
%
\newpage
\section{Introduction}

Employee attrition has long been recognized as a critical organizational challenge due to its direct economic costs, the loss of accumulated knowledge, and its negative effects on team stability and productivity \cite{ref_manafi2025,ref_gundlach2025}.

Addressing voluntary employee attrition requires both identifying employees who are at risk of leaving an organization and \textit{prioritising retention efforts toward those employees whose continued employment is most valuable within a given organizational context}.

In practice, organizations rarely have the capacity to apply retention strategies uniformly across all at-risk employees. Instead, rather than relying only on individual attrition risk, organisations should prioritise retention efforts based on, for example, organisation-specific combination of factors, such as strategic importance, difficulty of replacement, tenure, or accumulated experience.

The growing availability of human resources data has led to the widespread adoption of data-driven approaches for supporting attrition-related decisions.
In particular, machine learning techniques have been extensively used to model employee attrition as a supervised classification or risk scoring problem, estimating the probability that an employee will leave the organization based on demographic, job-related, and performance-related attributes \cite{ref_kambhampati2024,ref_biswas2023,ref_park2024,ref_habous2021,ref_jin2024}.
While these approaches have demonstrated strong predictive capabilities, they primarily focus on risk estimation rather than explicitly prioritizing retention.

More recent contributions incorporate explainability and interpretability mechanisms in order to support managerial decision-making and increase trust in predictive models \cite{ref_baydili2025,ref_konar2025}.
Although explainable models provide valuable insights into the drivers of attrition, they typically remain centered on a single analytical perspective and do not explicitly address how multiple assessments or organizational constraints should be combined to make decisions about retention measures.

Beyond automated prediction, alternative approaches rely on human-driven or expert-based processes, where attrition indicators are combined with managerial judgment, organizational policies, or qualitative evaluations to determine retention actions \cite{ref_mitravinda2022}.
Other works adopt decision support or multi-criteria perspectives, recognizing that attrition-related decisions involve multiple, potentially conflicting factors; however, these approaches often rely on a single aggregated analytical view or remain detached from heterogeneous automated assessments \cite{ref_manafi2025}.
Overall, existing approaches tend to conflate attrition management with attrition prediction, treating retention prioritization as a direct byproduct of risk estimation or as an informal managerial process.
The explicit modeling of attrition intervention as a cooperative decision-making problem, that is, by integrating multiple heterogeneous automated assessments and organizational context, remains largely unexplored.

The European Law Institute (ELI) defined automated decision making (ADM) system as ``\textit{a (computational) process, including AI techniques and approaches, that, fed by inputs and data received or collected from the environment, can generate, given a set of pre-defined objectives, outputs in a wide variety of forms (content, ratings, recommendations, decisions, predictions, etc.)}'' \cite{ref_eli2022_innovation}. Within this definition, ADM systems generate predictive outputs that support decision-making. The predictive outputs can be combined with any specific context for the purpose of prioritising employees retention. 
Context, in turn, refers to a stable layer of information that frames the problem under analysis. It is conceptually distinct from the problem itself, as decisions are always evaluated within a given situational, organizational, or operational setting that conditions their interpretation and relevance \cite{ref_dourish2004,ref_lamata2020}.

In this paper, we address this gap by proposing a cooperative and context-aware approach for prioritising employee retention.
Rather than aiming to improve predictive accuracy, the proposed framework structures heterogeneous automated assessments, originating from different organizational perspectives, into a transparent prioritization mechanism.
By explicitly modeling cooperation among decision components and context modulation, the approach supports informed retention-related decisions in settings characterized by distributed information, multiple stakeholders, and explicit resource constraints.

\section{Proposed Approach}

We conceptualize employee attrition intervention management as a cooperative, context-aware decision-making problem in which multiple heterogeneous ADM components contribute partial, complementary, and potentially conflicting assessments (Fig. \ref{fig:decision_flowchart}). 
\begin{figure}[t]
\centering
\includegraphics[width=1.0\textwidth, clip=true]{cadm.png}
\caption{Decision-making flowchart.}    
\label{fig:decision_flowchart}
\end{figure}
Rather than producing a single predictive output, the objective is to construct a structured prioritization mechanism that supports informed intervention decisions under organizational constraints.

The proposed approach builds on three core ideas.
First, attrition-related assessments are inherently distributed across \textit{organizational units} (OU) or departments, each operating with different information and objectives.
Second, effective intervention decisions require the \textit{integration} of these heterogeneous assessments in a principled and transparent manner.
Third, prioritization must be modulated by organizational \textit{context}, reflecting strategic considerations that extend beyond attrition risk alone.

Let $\mathcal{E} = \{e_1, e_2, \dots, e_n\}$ denote the set of employees under consideration.
Each employee $e_i$ is characterized by a feature vector $\mathbf{x}_i$, representing available demographic, professional, performance-related, and other context information.

The objective of this work is not to predict whether an employee $e_i$ will leave the company or not, but to construct a \textit{prioritization index} 
\[
\pi(e_i) \in \mathbb{R}
\]
, which induces an ordering over employees and enables decision-makers to identify a subset of $m \ll n$ employees for targeted retention interventions.

This prioritization index is derived from the cooperative integration of heterogeneous automated assessments and their modulation by organizational context.

We consider a collection of $K$ automated decision components
\[
\mathcal{A} = \{A_1, A_2, \dots, A_K\},
\]
each representing a distinct organizational unit or departmental perspective.
Each component operates on a subset of the available information and produces an independent assessment of attrition-related risk.

Formally, for an employee $e_i$, component $A_k$ outputs a scalar value
\[
p_k(e_i) \in [0,1],
\]
interpreted as the estimated likelihood of attrition according with organizational unit $k$.
These assessments are not assumed to be calibrated, mutually consistent, or jointly optimized.
Instead, their heterogeneity reflects the distributed nature of organizational knowledge and responsibility.

Not all decision components contribute equally reliable information.
Differences in data quality, modeling assumptions, and informational scope may affect their usefulness for prioritization. To account for this, each component $A_k$ is associated with a non-negative confidence weight $w_k$, satisfying $\sum_{k=1}^{K} w_k = 1$.

These weights encode the relative trust assigned to each component and can be derived from historical performance, validation metrics, expert judgment, or governance rules.
Their role is not to enforce dominance of a single component, but to modulate the cooperative integration process in a transparent and adjustable manner.

For each employee $e_i$, the partial assessments produced by the $K$ components are aggregated into an integrated risk indicator:
\begin{equation}\label{eq:colective_risk}
R(e_i) = \sum_{k=1}^{K} w_k \, p_k(e_i)    
\end{equation}

Attrition intervention decisions are rarely based on risk assessments alone.
Organizational priorities such as role criticality, replacement difficulty, accumulated firm-specific knowledge, or career trajectory considerations often play a decisive role. To incorporate such considerations, we introduce a set of $L$ factors that may influence the context $\mathcal{C} = \{C_1, C_2, \dots, C_L\}$. Each factor $C_\ell$ is modeled as a fuzzy proposition, allowing gradual and interpretable degrees of relevance rather than rigid thresholds \cite{ref_perezcanedo2023}.
For an employee $e_i$, the degree of satisfaction of such factor $C_\ell$ is denoted by $\mu_\ell(e_i) \in [0,1]$.

These membership values capture how well an employee aligns with organizational priorities associated with each context dimension.

The individual factors, that may influence the context, are combined into a collective context satisfaction score
\[
C(e_i) = H\big(\mu_1(e_i), \mu_2(e_i), \dots, \mu_L(e_i)\big),
\]
where $H: [0,1]^L\rightarrow[0, 1]$ denotes an aggregation operator.
In line with conservative organizational decision-making, conjunction-based operators (e.g., fuzzy AND) are particularly suitable, as they emphasize joint satisfaction of multiple criteria.

The final prioritization index is defined as a combination of cooperative risk and context relevance:
\[
\pi(e_i) = \Psi\big(R(e_i), C(e_i)\big),
\]
where $\Psi(\cdot)$ is a monotonic combination function. This definition ensures that employees are prioritized not only because they exhibit elevated attrition risk, but because intervening in their case aligns with organizational priorities and constraints.

It is important to remark that, by decoupling prediction, aggregation, and context modulation, the proposed approach reflects the structure of real-world organizational decision processes.
It does not seek to replace predictive analytics, but to embed them within a cooperative and interpretable prioritization framework that explicitly acknowledges heterogeneity, uncertainty, and organizational context.

\section{Illustrative Case Study}\label{sec:illustrative}

This section illustrates the proposed cooperative and context-aware prioritization approach using a public employee attrition dataset.
The objective is to demonstrate how heterogeneous automated assessments and factors that may influence the context interact to produce an interpretable prioritization of intervention candidates.

\subsection{Data Description}

The experiments are conducted using a publicly available Kaggle dataset, \textit{IBM HR Analytics Employee Attrition and Performance}
\footnote{\url{https://www.kaggle.com/datasets/pavansubhasht/ibm-hr-analytics-attrition-dataset} (Checked on 19/03/2026)}, which consists of 1.470 employee records and 35 attributes. The attributes involve demographic, professional, performance-related and organizational variables, as well as a binary attrition indicator. The dataset is split into training (80\%) and test (20\%) partitions. The test set is used solely to estimate the relative reliability of the decision components.

\subsection{Instantiation of Heterogeneous Decision Components}

Following the formulation in Section~2, four automated decision components are instantiated, each representing a distinct organizational perspective.
Each component was implemented as a binary classification model trained using an AutoML pipeline based on the \textit{H2O} framework \cite{ref_h2o2026}. Models were trained using 5-fold cross-validation and a Bernoulli distribution, consistent with a binary classification setting, yielding probabilistic risk estimates suitable for downstream aggregation.

To reflect organizational heterogeneity and data privacy, each component operates on a subset of employee attributes, corresponding to the type of information typically available to a given organizational unit.
Table~\ref{tab:departments} summarizes the organizational unit perspectives and their representative feature subsets. Each decision component $A_k$ outputs, for an employee $e_i$, an attrition risk score.

\begin{table}[t]
\centering
\setlength{\tabcolsep}{6pt}
\caption{Organizational unit perspectives and representative feature subsets} 
\label{tab:departments}
\renewcommand{\arraystretch}{1.3}
\begin{tabular}{p{2.4cm}|p{9cm}}
\hline
\textbf{Organizational unit} & \textbf{Representative attributes (features)} \\
\hline
Culture and Wellbeing ($A_{1}$) &
\textit{Age, Gender, Department, JobRole, JobLevel, EnvironmentSatisfaction, JobInvolvement, JobSatisfaction, RelationshipSatisfaction, WorkLifeBalance, OverTime, BusinessTravel} \\
\hline
Finance ($A_{2}$) &
\textit{Age, Gender, Department, JobRole, JobLevel, MonthlyIncome, MonthlyRate, DailyRate, HourlyRate, PercentSalaryHike, StockOptionLevel} \\
\hline
Operations and Development ($A_{3}$) &
\textit{Age, Gender, Department, JobRole, JobLevel, TotalWorkingYears, NumCompaniesWorked, YearsAtCompany, YearsInCurrentRole, YearsSinceLastPromotion, YearsWithCurrManager, PerformanceRating, TrainingTimesLastYear} \\
\hline
Human Resources ($A_{4}$) &
\textit{Age, Gender, Department, JobRole, JobLevel, MaritalStatus, Over18, Education, EducationField, DistanceFromHome, StandardHours} \\
\hline
\end{tabular}
\end{table}

After initializing the AutoML procedure, three models, corresponding to the best, average, and lowest performance, were selected for each subset. Model selection was based on the log-loss criterion. For the final attrition risk estimation, we prioritized models with average performance within four organizational units (Table 2). This decision was based on the need to balance predictive accuracy with generalizability, avoiding overfitting to extreme cases. For each subset, the attrition risk was computed as a weighted average of the predicted probabilities of positive attrition, where weights correspond to the relative contribution of each model in the ensemble. The relative confidence of each decision component is quantified using the Area Under the ROC Curve (AUC) computed on the test partition, as shown in Table \ref{tab:ou_models}, and reflect the reliability of each model in distinguishing between attrition and non-attrition cases.

\begin{table}[h!]
\centering
\caption{Models selected for the ensemble, expressed as parametrized functions, with corresponding AUC scores and normalized weights.}
\setlength{\tabcolsep}{8pt}
\renewcommand{\arraystretch}{1.3}
\begin{tabular}{|l |l| c |c|}
\hline
\textbf{OU} & \textbf{Model} & \textbf{AUC} & \textbf{Weight} \\
\hline
A1 & GBM(\textit{NT}=247, \textit{MTD}=17, \textit{FOL}=15) & 0.99 & 0.35 \\
\hline
A2 & GBM(\textit{NT}=196, \textit{MTD}=13, \textit{FOL}=30) & 0.98 & 0.30 \\
\hline
A3 & GBM(\textit{NT}=159, \textit{MTD}=7,  \textit{FOL}=15) & 0.95 & 0.25 \\
\hline
A4 & DL(\textit{E}=8, \textit{IDR}=0.15)                   & 0.73 & 0.10 \\
\hline
\end{tabular}

\vspace{1.5mm}
\footnotesize{
\textit{NT}: number of trees; 
\textit{MTD}: max tree depth; 
\textit{FOL}: min number of (weighted) observations in a leaf; 
\textit{E}: number of training epochs; 
\textit{IDR}: input-layer dropout ratio.
}
\label{tab:ou_models}
\end{table}
\subsection{Context definition} 
In this section, we considered two independent contexts that may influence companies' decisions: 1) one related with an economic crisis; 2) and another related with digital or technological transformation.

Each context is modeled using a set of fuzzy propositions encoding organizational priorities and constraints. A fuzzy proposition is a statement $p$ whose falsity or truth is expressed to varying degrees within the interval $[0, 1]$ and is denoted by $\text{Truth}(p):\;V \text{ is } \tilde{F}$, where $V$ is variable that takes values $x$ in a universe of discourse $X$ and $\tilde F$ is a fuzzy set on $X$ that represents a fuzzy predicate (linguistic term). The degree of truth of $p$ is the same as the degree of membership of $x$ in $\tilde F$,  i.e., $\Truth(p)=\mu_{\tilde F}(x)$.
\subsubsection*{\textit{Context 1: }Economic Crisis.}
\label{app:economic_crisis}
This context represents a scenario of financial pressure, in which companies are forced to make decisions about retaining staff in conditions of limited resources and heightened uncertainty. In such situations, the main objective is not to maximise long-term growth, but to ensure the short- and medium-term stability of the company while controlling operating costs.

Specifically, this context prioritizes the retention of employees who satisfy, to varying degrees, the following set of fuzzy propositions:
\begin{itemize}
    \item \textit{Organizational tenure is medium};
    \item \textit{Performance is high}; 
    \item \textit{Salary cost (within the department) is average};
    \item \textit{Commitment to work is satisfactory}.
\end{itemize}

Unless otherwise stated, all the linguistic terms are defined with a triangular membership function with three parameters $(a,b,c)$:
\[
\mu_A(x) = \begin{cases} 
      0 & x < a \\
      (x-a)/(b-a) & a \leq x < b\\
      (c-x)/(c-b) & b \leq x \leq c \\
      0 & x > c, 
   \end{cases}
\]
where x is the numerical value of the criterion being evaluated for a given employee.   

The definitions for these criteria are described below.

\paragraph{Organizational Tenure is medium.} This proposition is based on the variable $YearsAtCompany$. Let $x = \text{YearsAtCompany}$; in the dataset under study, $x \in [0,36]$. It is used as an indicator of work experience and operational reliability.

The linguistic term \textit{medium} for a given number of years has the parameters $(a=5, b = 10, c=15)$, where the central value (b = 10) represents a typical mid-career point in the dataset, and the bounds (a, c) define a symmetric range around this value corresponding to employees with moderately low and moderately high tenure.

In times of economic crisis, it is assumed that employees with average period of tenure provide an optimal balance between accumulated organisational knowledge and economic efficiency. 

\paragraph{Performance is at least medium.} Performance reflects an employee's ability to meet the expectations associated with their position and contribute to the achievement of the company's goals. In times of economic crisis, the company seeks to retain employees whose performance is at least medium.

The linguistic term \textit{at least medium} for the variable Performance is defined over the variable $x = \text{PerformanceRating}$; in the dataset under study, with $x \in \{1,2,3,4,5\}$ (the discrete 1-5-point performance rating scale) as: 

\[
\mu_{\text{PerfRating}}(x) =
\begin{cases}
0, & x \le 1, \\
\dfrac{x - 1}{3}, & 1 < x < 4, \\
1, & x \ge 4.
\end{cases}
\]

Ratings of 4 and 5 correspond to clearly satisfactory performance, and a linear transition between 2 and 3 provides a gradual increase in membership for employees approaching the “average” threshold.

\paragraph{Salary cost is average.} In the context of an economic crisis, salary cost is a critical factor in employee retention decisions. However, absolute salary values cannot be directly compared between departments due to structural differences in pay scales and job requirements. As shown in Table \ref{tab:salary}, the average and maximum salary levels in the human resources (HR) department are significantly lower than in the research and development (R\&D) and sales (SAL) departments.
\vspace{-1.5em}
\begin{table}
\centering
\small 
\caption{Salary ranges and averages by department\label{tab:salary}}
\begin{tabular}{|l|c|c|c|}
\hline
Department & Min & Max & Average \\ \hline
Human Resources        & 1.555 & 10.482 & 6.019 \\ \hline
Research \& Development & 1.009 & 19.740 & 10.375 \\ \hline
Sales                  & 1.081 & 19.845 & 10.463 \\ \hline
\end{tabular} 
\end{table}
\vspace{-1.5em}

The linguistic term \textit{average} for the variable Salary is defined with a triangular membership function, and the corresponding parameters for every department are: 

\begin{center}
\small
\begin{tabular}{|l|c|c|c|}
\hline
Param. & a & b & c \\
\hline
$\mu_{\text{HR-avg}}$   & 3000 & 6000 & 8000  \\
\hline
$\mu_{\text{RD-avg}}$    & 7000 & 10000 & 14000 \\
\hline
$\mu_{\text{Sales-avg}}$ &  7000 & 10000 & 14000 \\
\hline
\end{tabular} 
\end{center}

These parameters are derived from the empirical salary distribution within each department. The central value $b$ is set close to the observed average salary (Table \ref{tab:salary}), while $a$ and $c$ define a symmetric range around this mean, capturing moderately low and moderately high salaries relative to the department.

\paragraph{Commitment to work.}
Commitment to work depends on two variables: job involvement and job satisfaction, that are measured on the same ordinal scale, with higher values reflecting greater engagement or job satisfaction.

Job involvement  reflects behavioural and motivational engagement, while job satisfaction reflects emotional evaluation of work. 

The context requires that:
\begin{center}
    \texttt{JobInvolvement is satisfactory AND JobSatisfaction is satisfactory}
\end{center}

Let $x = \text{JobInvolvement}$ and $y = \text{JobSatisfaction}$ with $x,y \in \{1,2,3,4\}$. 
The different degrees of membership to the linguistic term \textit{Satisfactory} are:
\begin{center}
\begin{tabular}{|l|c|c|c|c|}
\hline
Value & 1 & 2 & 3 & 4 \\
\hline
$\mu_{\text{invol}}$ & 0.15 & 0.50 & 0.85 & 1.0 \\
\hline
$\mu_{\text{satis}}$ & 0.10 & 0.40 & 0.70 & 1.0 \\
\hline
\end{tabular}
\end{center}

A min operator was used to combine the two variables, ensuring that overall commitment reflects the lowest level of satisfaction or engagement, consistent with standard fuzzy logic practice.
\subsubsection*{\textit{Context 2:} Digital or Technological Transformation.}
\label{app:digital_transformation}

This context captures a strategic scenario oriented toward future capabilities rather than short-term cost containment. Focus on automation of processes, adoption of AI or data-driven decision-making, new digital tools.
Retention is prioritized for employees with strong upskilling potential, adaptable roles, scalable career paths,
and an adequate investment horizon.

This context is oriented toward retaining employees who satisfy the following set of fuzzy propositions:
\begin{itemize}
    \item \textit{Upskilling capacity is at least medium or high}; 
    \item \textit{Role adaptability is high};
        \item \textit{Career scalability is at least medium or higher}.
\end{itemize}

The criteria used to define this context are described below.
\paragraph{Upskilling capacity is at least medium or high.} 
The proposition is based on the variable $x = \text{TrainingTimesLastYear}$ which, in the dataset under study, with $x \in \{0,1,2,3,4,5,6\}$.

In this context, the company prioritises employees who demonstrate the ability and willingness to get new skills. Participation in training activities is seen as a measurable indicator of interest in learning and openness to continuous professional development.
\begin{table}[h]
\centering
\small
\begin{tabular}{|l|c|c|c|c|c|c|c|c|}
\hline
TrainingTimesLastYear & 0 & 1 & 2 & 3 & 4 & 5 & 6 \\
\hline
$\mu_{\text{upskil}}$ & 0.00 & 0.20 & 0.30 & 0.50 & 0.85 & 0.95 & 1.00 \\
\hline
\end{tabular}
\end{table}
\vspace{-2.4em}

\paragraph{Role adaptability.} An employee is considered adaptable if their education and current position indicate experience in various fields of expertise or non-routine work. Role adaptability is assessed by analyzing the match between the employee's educational background and their position, with greater weight given to non-standard (cross-industry) career paths, taking into account the adaptability requirements of the given position.

Two categorical variables are used: 1. EducationField, representing formal educational background; 2. JobRole, representing the functional position within the organization.

Certain education-role combinations are considered standard or expected, and therefore indicative of specialization rather than adaptability. These combinations are assigned a low but non-zero relevance value. The following matches are considered standard: 

\begin{enumerate}
    \item Marketing $\leftrightarrow$ Sales Executive, Sales Representative 
    \item Medical / Life Sciences / Technical Degree $\leftrightarrow$ Research Director, Manufacturing Director, Research Scientist, Laboratory Technician 
    \item Human Resources $\leftrightarrow$ Human Resources 
\end{enumerate}

For all other education-role combinations, adaptability relevance is determined by the internal adaptability requirements of the given position, regardless of the field of education. These values reflect expert assessments of the cognitive complexity, decision-making independence, and transferability of skills associated with each role. 

The degrees of membership in the fuzzy set of role adaptability are based on expert judgment and represent relative adaptability rather than absolute skill levels, and presented below.

\begin{table}[h]
\centering
\label{tab:role_adaptability}
\begin{tabular}{|l|c c|l|c|}
\hline
JobRole & $\mu_{\text{role}}$ &\hspace{1.3cm}& JobRole & $\mu_{\text{role}}$\\
\hline
Research Director & 1.00 && Manufacturing Director & 0.95 \\
\hline
Manager & 0.70 && Research Scientist & 0.85 \\
\hline
Sales Executive & 0.70 && Healthcare Representative & 0.80 \\
\hline
Laboratory Technician & 0.80 && Human Resources & 0.50 \\
\hline
Sales Representative & 0.60 && & \\
\hline
\end{tabular}
\end{table}


Let \(E = \text{EducationField}\) and \(R = \text{JobRole}\).  
The role adaptability relevance \(\mu_{\text{adaptability}}\) is defined as:
\[
\mu_{\text{adapt}} =
\begin{cases}
0.2, & \text{if } (E,R) \text{ corresponds to a standard education--role match} \\
\mu_{\text{role}}(R), & \text{otherwise}
\end{cases}
\]

This rule ensures that specialization paths are not excluded but weakly prioritized and non-standard career paths are rewarded in proportion to the adaptability demands of the role. Adaptability remains interpretable and bounded in \([0,1]\).

\paragraph{Career scalability}is at least medium. Career scalability $cs$ is measure as\\
\vspace{-1.5em}
\begin{center} 
    \text{x = 1 - (YearsInCurrentRole/YearsAtCompany)}
\end{center}
    
If an employee remains in the same position for many years and does not receive any promotion, this can also be interpreted as career stagnation.

The linguistic term \emph{at least medium} for Career Stability is defined as:

\[
\mu_{\text{CS}}(x) =
\begin{cases}
\dfrac{x}{0.5}, & 0 \leq x < 0.5 \\[6pt]
1, & x \geq 0.5
\end{cases}
\] 


\subsection{Prioritization Procedure}

For each selected employee $e_i$, the partial risk assessments produced by the decision components are aggregated into a collective risk score as defined in Eq. \ref{eq:colective_risk}. In parallel, the contextual membership values $\mu_\ell(e_i)$ are computed and combined using a conjunction-based aggregation operator to obtain a collective context satisfaction score $C(e_i)$. Finally, we defined the prioritization index
as 
\begin{equation}
\pi(e_i) = \lambda\, R(e_i) + (1 - \lambda)\, C(e_i)
\end{equation}
where $\lambda \in [0,1]$ is a policy-level weighting factor defined by the organization, reflecting the relative importance assigned to attrition risk versus context relevance in the prioritization process.

\subsection{Results}

To assess attrition risk in a contextual framework, we combined the positive attrition risk scores with context-relative scores, as illustrated in Figure \ref{fig:risk_attrition}.

\begin{figure}[t]
    \centering
    \begin{subfigure}{0.48\textwidth}
        \centering
        \includegraphics[width=\textwidth]{Context 1}
        \caption{Context 1}
        \label{fig:context1}
    \end{subfigure}
    \hfill
    \begin{subfigure}{0.48\textwidth}
        \centering
        \includegraphics[width=\textwidth]{Context 2}
        \caption{Context 2}
        \label{fig:Context 2}
    \end{subfigure}
    \caption{Risk attrition across contexts}
    \label{fig:risk_attrition}
\end{figure}


Attrition risk scores were estimated in the interval $[0,1]$ and combined with context scores within the same interval $[0,1]$. A threshold of $R(e_i) \geq 0.8$ was adopted to identify employees at high risk of attrition. A threshold of $C(e_i) \geq 0.5$ was further applied to capture cases with moderate to high contextual vulnerability. Critical cases were defined as those located in the top-right quadrant of the risk-context plane, i.e., employees simultaneously satisfying $R(e_i) \geq 0.8$ and $C(e_i) \geq 0.5$.

For Context 1, three employees met the criteria, and all were retained for analysis. For Context 2, twenty-two employees met the criteria; to illustrate the most critical cases, the three employees with the highest context scores were selected. For each of these six employees, the prioritization index $\pi(e_i)$ was calculated with $\lambda = 0.5$. After that we evaluated the distribution of scores across contexts.

\paragraph{Context 1: Critical Case Analysis.}
The individual scores of the employees selected for analysis who met the dual threshold condition of high attrition risk and moderate-to-high contextual vulnerability are summarized below.
\begin{table}[h!]
\centering
\begin{tabular}{|c|c|c|c|c|}
\hline
\textbf{Employee} & \textbf{Attrition Risk} & \textbf{Context Score} & \textbf{Priority Index} & \textbf{Attrition} \\
\textbf{ID} & $R(e_i)$ & $C(e_i)$ & $\pi(e_i)$ &  \textbf{original} \\
\hline
1421 & 0.90 & 0.51 & 0.71 & Yes \\
\hline
118  & 0.88 & 0.50 & 0.69 & Yes \\
\hline
1457 & 0.84 & 0.50 & 0.67 & Yes \\
\hline
\end{tabular}
\end{table}
\vspace{-1.5em}

These cases exhibit very high attrition risk (>=0.84), with context scores clustered around 0.50-0.51, just above the vulnerability threshold.

The resulting prioritization scores range from 0.67 to 0.71, indicating that these cases are risk-dominant: the attrition risk contributes more heavily to $\pi(e_i)$ than context.

This means that in Context 1, the prioritization index is risk-dominant: employees are flagged as critical primarily because of their high probability of leaving, rather than because of fragility of context. The context dimension plays only a minimal role in distinguishing between them.

\paragraph{Context 2: Critical Case Analysis.}
Three employees were selected from the subset satisfying the dual-threshold condition of high attrition risk ($R(e_i) \geq 0.8$) and moderate-to-high vulnerability threshold ($C(e_i) \geq 0.5$).Their individual scores are summarized below.

\begin{table}[h!]
\centering
\begin{tabular}{|c|c|c|c|c|}
\hline
\textbf{Employee} & \textbf{Attrition Risk } & \textbf{Context Score } & \textbf{Priority Index} & \textbf{Attrition} \\
\textbf{ID} & $R(e_i)$ & $C(e_i)$ & $\pi(e_i)$ &  \textbf{original} \\

\hline
297   & 0.88 & 0.85 & 0.86 & Yes\\
\hline
1767  & 0.88 & 0.85 & 0.86 & Yes\\
\hline
1572  & 0.87 & 0.85 & 0.86 & Yes\\
\hline
\end{tabular}
\end{table}


All three employees exhibit very high attrition risk (>=0.87) and consistently high context scores (0.85), well above the vulnerability threshold.

The resulting prioritization scores are uniformly high (0.86), indicating that these cases are jointly driven by both attrition risk and fragility of context.

Unlike Context 1, where context scores were minimal and risk dominated, Context 2 shows strong alignment between risk and context, suggesting that these employees are not only likely to leave but are embedded in environments that may be contributing to that risk.

The presented results highlight employees with the highest perceived attrition risk and context vulnerability, demonstrating the proposed framework's ability to differentiate cases based on risk and context. While we cannot fully confirm actual attrition cases, this approach provides a structured method for prioritizing interventions in resource-constrained settings.


\section{Conclusions}

The comparative analysis of the results for two different contexts highlight distinct dynamics in the drivers of employee attrition risk. In Context 1, the narrow spread of context scores indicates limited contextual differentiation, suggesting that attrition in this group is risk-dominant and interventions should focus on direct retention measures.
In contrast, Context 2 revealed employees with high attrition risk and consistently elevated context scores. The prioritization indices for these cases reflects a balanced contribution of risk and context. This pattern demonstrates that attrition in Context 2 is jointly driven by individual probability and structural vulnerability. Consequently, effective interventions must combine retention strategies with systemic improvements to the work environment.

Importantly, these findings emphasize that decision-making cannot rely solely on machine learning outputs. While predictive models provide valuable probabilities of attrition, the contextual dimension ($C(e_i)$) adds critical interpretive power. Context scores highlight structural or environmental conditions that shape employee behavior, guiding managers toward interventions that are not only statistically justified but also organizationally meaningful. Thus, attrition analysis must integrate both quantitative risk estimates and qualitative contextual insights to ensure that decisions are robust, actionable, and aligned with organizational realities.

\begin{credits}
\subsubsection*{\ackname} Mariia Godz has a pre-doctoral research contract associated with the scientific project “Study, Analysis and Evaluation of Cooperative Automated Decision-Making Systems (CADEMAS)” with reference number PID2023-146575NB-I00, funded by MCIU/AEI/I0.13039/501100011033  including Feder Funds/UE, and by FSE+.

\end{credits}

\begin{thebibliography}{99}
\bibitem{ref_baydili2025}
Baydili, \. {I}.T., Tasci, B.: Predicting employee attrition: XAI-powered models for managerial decision-making. Systems \textbf{13}(7), 583 (2025)

\bibitem{ref_biswas2023}
Biswas, A.K., Seethalakshmi, R., Mariappan, P., Bhattacharjee, D.: An ensemble learning model for predicting the intention to quit among employees using classification algorithms. Decision Analytics Journal \textbf{9}, (2023)

\bibitem{ref_dourish2004}
Dourish, P.: What we talk about when we talk about context. Personal and Ubiquitous Computing \textbf{8}(1), 19--30 (2004). \doi{10.1007/s00779-003-0253-8}

\bibitem{ref_eli2022_innovation}
ELI Innovation Paper: Guiding principles for automated decision-making in the EU. European Law Institute (2022). Available at: \href{https://www.europeanlawinstitute.eu/fileadmin/user_upload/p_eli/Publications/ELI_Innovation_Paper_on_Guiding_Principles_for_ADM_in_the_EU.pdf}{ELI Innovation Paper}, last accessed 2026/01/19.


\bibitem{ref_gundlach2025}
Gundlach, H.A.D.: What really influences teacher attrition, migration, and retention? The Australian Educational Researcher \textbf{52}(5), 3079--3099 (2025). \doi{10.1007/s13384-025-00842-4}

\bibitem{ref_h2o2026}
H2O.ai Homepage, \url{https://www.h2o.ai/}, last accessed 2026/01/10.

\bibitem{ref_habous2021}
Habous, A., Nfaoui, E.H., Oubenaalla, Y.: Predicting employee attrition using supervised learning classification models. In: 2021 Fifth International Conference on Intelligent Computing in Data Sciences (ICDS), pp.~1--5. IEEE (2021). \doi{10.1109/ICDS53782.2021.9626761}

\bibitem{ref_jin2024}
Jin, A., Yang, Z., Liu, X., Lou, S., Zhuang, Z., Zhang, S., Zhang, C.: Predicting employee attrition using machine learning approaches. In: 2024 9th International Symposium on Computer and Information Processing Technology (ISCIPT), pp.~419--426. IEEE (2024). \doi{10.1109/ISCIPT61983.2024.10672901}

\bibitem{ref_kambhampati2024}
Kambhampati, V.K., Rao, K.B.V.B.: Advancing employee attrition models: A systematic review of machine learning techniques and emerging research opportunities. In: 2024 8th International Conference on I-SMAC (IoT in Social, Mobile, Analytics and Cloud) (I-SMAC), pp.~1100--1107. IEEE (2024). \doi{10.1109/I-SMAC61858.2024.10714595}

\bibitem{ref_konar2025}
Konar, K., Das, S., Das, S., Misra, S.: Employee attrition prediction using Bayesian optimized stacked ensemble learning and explainable AI. SN Computer Science \textbf{6}(6) (2025). \doi{10.1007/s42979-025-04204-w}

\bibitem{ref_lamata2020}
Lamata, M.T., Pelta, D.A., Verdegay, J.L.: The role of the context in decision and optimization problems. In: Fuzzy Approaches for Soft Computing and Approximate Reasoning: Theories and Applications, pp.~75--84. Springer, Cham (2020). \doi{10.1007/978-3-030-54341-9_7}

\bibitem{ref_manafi2025}
Manafi Varkiani, S., Pattarin, F., Fabbri, T., Fantoni, G.: Predicting employee attrition and explaining its determinants. Expert Systems with Applications \textbf{272}, (2025). \doi{10.1016/j.eswa.2025.126575}

\bibitem{ref_mitravinda2022}
Mitravinda, K.M., Shetty, S.: Employee attrition: Prediction, analysis of contributory factors and recommendations for employee retention. In: 2022 IEEE International Conference for Women in Innovation, Technology \& Entrepreneurship (ICWITE), pp.~1--6. IEEE (2022). \doi{10.1109/ICWITE57052.2022.10176235}

\bibitem{ref_park2024}
Park, J., Feng, Y., Jeong, S.-P.: Developing an advanced prediction model for new employee turnover intention utilizing machine learning techniques. Scientific Reports \textbf{14}(1) (2024). \doi{10.1038/s41598-023-50593-4}

\bibitem{ref_perezcanedo2023}
Pérez-Cañedo, B., Porras, C., Pelta, D.A., Verdegay, J.L.: Modeling contexts as fuzzy propositions in optimization problems. IEEE Transactions on Fuzzy Systems \textbf{31}(5), 1474--1483 (2023). \doi{10.1109/TFUZZ.2022.3203786}
\end{thebibliography}
\appendix

\end{document}


